% =============================================================================
\section{Tarefa 6 --- Gestos e estimação da trajetória}
\label{sec:gestos}
% =============================================================================

\subsection{Da orientação à posição}

Com a orientação $R(t)$ da Tarefa~5, a aceleração do sensor em $\{w\}$
obtém-se de~\eqref{eq:acc}, e a posição integrando duas vezes:
\begin{equation}
\vect a = R\,S_a^{-1}(\med a - \vect b_a) + \vect g, \qquad
\vect v(t) = \int_0^t \vect a\,d\tau, \qquad
\vect p(t) = \int_0^t \vect v\,d\tau .
\label{eq:pos}
\end{equation}

\begin{destaque}
\textbf{Previsão.} Um erro constante $b$ na aceleração dá um erro de posição
$\tfrac12\,b\,t^2$. Há duas fontes: o desvio residual do acelerómetro, e o
erro de inclinação $\delta\theta$, que deixa por subtrair $g\sin\delta\theta$.
Com \ang{1} de erro de inclinação são \SI{0.17}{\meter\per\second\squared},
ou \SI{77}{\centi\meter} ao fim de \SI{3}{\second}. Por isso os gestos são
curtos e começam e acabam com o sensor parado.
\end{destaque}

Como cada gesto começa e acaba em repouso, a velocidade nos dois extremos é
conhecida (zero). \todo{correção usada: retirar à velocidade a deriva linear
entre os dois instantes de repouso; e, tudo sendo pós-processado, um
suavizador~\cite{rauch1965} em vez de um filtro só para a frente.}

\subsection{Validação com o robô (ensaios F e G)}

Antes dos gestos à mão, o método foi aplicado a dois movimentos executados
pelo robô, de que se conhece a posição em cada instante: um ``acenar'' e um
``soco'', desenhados nas juntas como o swing do Laboratório~2.

\begin{figure}[H]
\centering
\fig{Photos/deriva.png}{0.75\textwidth}
\caption{\todo{erro de posição ao longo do tempo, com o sensor parado e numa
trajetória do robô, com a curva prevista $\tfrac12 b t^2$; com e sem
calibração e com e sem correção de velocidade nula.}}
\label{fig:deriva}
\end{figure}

\todo{erro de posição no fim de cada gesto do robô; até que duração o erro
fica abaixo de um valor útil (por exemplo \SI{5}{\centi\meter}).}

\subsection{Gestos à mão}

Com o melhor $w$ da Tarefa~5 e o sensor seguro na mão, fizeram-se devagar um
gesto de acenar e um de soco, \todo{$n$} vezes cada.

\begin{figure}[H]
\centering
\fig{Photos/gestos.png}{\textwidth}
\caption{\todo{acenar e soco à mão: orientação ao longo do tempo e trajetória
estimada, com fotogramas do vídeo; referência geométrica (comprimento da
régua, raio do antebraço).}}
\label{fig:gestos}
\end{figure}

\todo{resultados: amplitude e forma de cada gesto contra a referência;
repetibilidade entre repetições; o que se reconhece e o que não.}
